\documentclass[10pt]{beamer}
% Theme
\usetheme{metropolis}
% Get colors
\usepackage{xcolor}
% Get source code syntax highlighting
\usepackage[cache=false]{minted}
% Get CC doc license
\usepackage[type={CC},modifier={by-nc-sa},version={3.0}]{doclicense}

% Document information
\newcommand*{\mytitleTitle}{C++ Compilation and Pointers}
\newcommand*{\mytitleSubtitle}{StuPro Cheat Slides}
\newcommand*{\mytitle}{\mytitleSubtitle: \mytitleTitlet}
\newcommand*{\myauthor}{Niklas Mikeler}
\newcommand*{\mysubject}{StuPro - Multi-Kamera-System-Panorma-Aufnahmen - SS19}
\newcommand*{\mydate}{16.04.2019}

% Custom colors
\definecolor{lightGreyCustom}{HTML}{e8e8e8}
\definecolor{greyCustom}{HTML}{929292}

% Document methods
\newcommand{\urlCustom}[1]{\href{https://#1}{\textcolor{greyCustom}{\mintinline{sh}{#1}}}}

% Title page
\title{\mytitleTitle}
\subtitle{\mytitleSubtitle}
\date{\mydate}
\subject{\mysubject}
\keywords{\mytitleTitle, \mytitleSubtitle, \mysubject}
\author{\myauthor}
\institute{Universit{\"a}t Stuttgart}
\titlegraphic{\hfill\includegraphics[height=1.5cm]{resources/picture_unistuttgart_logo_deutsch.pdf}}

% Footer
\setbeamertemplate{frame footer}{\mytitleSubtitle: \mytitleTitle \ by \myauthor}


\begin{document}

\maketitle

\begin{frame}{Table of contents}
	\setbeamertemplate{section in toc}[sections numbered]
	\tableofcontents[hideallsubsections]
\end{frame}

\section{C++}

\begin{frame}{C++ $\gg$ What is a computer program?}
\begin{itemize}
	\item A computer program (aka. app/application) is a set of instructions that the computer can execute to do something
	\item The computer natively can only understand one language which is called machine code (\mintinline{sh}{10110000 01100001}) [Of course there could be a program in machine code that takes some text and runs commands based on this $\rightarrow$ Interpreter]
	\item The instruction set though differs from CPU to CPU (computer to computer) which means a native program needs to be rewritten for each new CPU, platform, etc.
\end{itemize}
\end{frame}

\begin{frame}{C++ $\gg$ Assembly?}
\begin{itemize}
	\item Because machine code is really difficult to read/write the assembly language was invented
	\item Instructions are now quite a bit more readable (\mintinline{as}{mov al, 061h}) and an additional software called the assembler takes the written code and translates it into machine code so that the CPU can understand it
	\item Because assembly code is very near to the hardware (\textit{low-level}) of each program code written in it is still not simply portable to another computer without doing major changes
\end{itemize}
\end{frame}

\begin{frame}{C++ $\gg$ High-level programming languages?}
\begin{itemize}
	\item To make source code (the code that the programmer should write) much more readable (\mintinline{cpp}{a = 97;}) and even portable (execute the program on any computer) a new type of programming languages were created
	\item These are called \textit{high-level} because they are designed to abstract the hardware to make programs in most cases completely independent from any hardware
	\item There are two ways a high-level programming language program can be run on a computer
\end{itemize}
\end{frame}

\begin{frame}{C++ $\gg$ Compiler and interpreters?}
\begin{enumerate}
	\item Compiler (C++, C, Haskell, Pascal)
	\begin{itemize}
		\item A compiler is a program that reads source code and produces a stand alone machine code executable and optimizes the code along the way and translates in rare cases even better to machine code than a human with assembly code would do
		\item This means that if the compiler supports translating to a specific CPU/instruction set/etc. any code written in this language supports this platform now (at least it should!)
	\end{itemize}
	\item Interpreter (Java, JavaScript, Python)
	\begin{itemize}
		\item An interpreter is a program that can run instructions from any input without it being necessarily in machine code
		\item This means that on any platform where an interpreter for the language exists can (should) automatically run these instructions
		\item Though the interpreter needs to read the source code each time the program gets executed and is because of the translation process each time often much more flexible on the language side (old programs run still/faster/more stable than before), but generally also a bit slower
	\end{itemize}
\end{enumerate}
\end{frame}

\begin{frame}{C++ $\gg$ What was its predecessor?}
\begin{itemize}
	\item Before C++ existed there was C which was developed in 1972 by Dennis Ritchie which allowed programmers to write imperative (procedural) programs
	\item It was standardized in 1989 and is till today widely used (C99 in 1999, C11 in 2011)
	\item It was designed to allow low-level access to memory but still have portability
\end{itemize}
\end{frame}

\begin{frame}{C++ $\gg$ Why does it exist?}
\begin{itemize}
	\item Bjarne Stroustrup developed C++ as an extension of the C language also known as \textit{C with Classes}
	\item It extends Cs imperative design by object-oriented (objects) and generic programming (templates) features while still providing low-level memory manipulation
	\item Like C C++ also has next to more high-level features still a design bias based on performance/efficiency/flexibility
	\item Since C++11 there are general-purpose smart pointers
	\item Since C++14 there is a better template and lambda support
	\item C++17 is since December 2017 the latest standardized version
	\item Like C it has the philosophy \textit{Trust the programmer} which means the programmer can make really fast and great programs but also the language doesn't stop him/her from making dangerous decisions 
\end{itemize}
\end{frame}

\section{C++ Compilation}

\begin{frame}{C++ Compilation $\gg$ A general overview}
\begin{enumerate}
	\item We have many C/C++ (\mintinline{sh}{.c}/\mintinline{sh}{.cpp}) source files
	\item Each source file (\mintinline{sh}{.c}/\mintinline{sh}{.cpp}) will be compiled by an independent compiler (This is cool, because files that haven't changed do not necessarily need to compiled again)
	\begin{enumerate}
		\item The preprocessor (preprocessor statement: \#\mintinline{sh}{xxx}) runs over the current file
		\begin{itemize}
			\item This will for example copy each mentioned header file and its parents directly in the document
			\inputminted[bgcolor=lightGreyCustom,fontsize=\scriptsize]{sh}{./resources/preprocessor_statement.h}
		\end{itemize}
		\item The file then gets compiled and transformed into an object (\mintinline{sh}{.o})
	\end{enumerate}
	\item Then each object is linked into one executable (more to the linking process later)
\end{enumerate}
\end{frame}

\begin{frame}{C++ Compilation $\gg$ A trivial C++ example}
\begin{enumerate}
	\item We can build it with the following commands:
	\inputminted[bgcolor=lightGreyCustom,fontsize=\scriptsize,firstline=3,lastline=4]{sh}{./resources/Compilation_example/build.sh}
\end{enumerate}
\vspace{-5mm}
\begin{columns}
	\begin{column}{0.57\textwidth}
		\begin{enumerate}\setcounter{enumi}{1}
			\item \mintinline{sh}{Main.cpp} (contains main function)
			\inputminted[bgcolor=lightGreyCustom,fontsize=\scriptsize]{c}{./resources/Compilation_example/Main.cpp}
		\end{enumerate}
	\end{column}
	\begin{column}{0.57\textwidth}
		\begin{enumerate}\setcounter{enumi}{2}
			\item \mintinline{sh}{Log.cpp}
			\inputminted[bgcolor=lightGreyCustom,fontsize=\scriptsize]{c}{./resources/Compilation_example/Log.cpp}
		\end{enumerate}
	\end{column}
\end{columns}
\end{frame}

\begin{frame}{C++ Compilation $\gg$ Why do we need header files?}
\begin{itemize}
	\item When the main function wants to call the function \mintinline{c}{add(5,6)}, it needs to know what it returns and what it takes in as argument before it can call the function, this information can come from a local declaration( \mintinline{c}{int add(int a, int b);})/definition or from a included header
	\item $\Rightarrow$ The compiler for each source file needs to know the size of structures and function arguments given through the function/symbol declaration, but not necessarily their definition
	\item This is why you include header files - else you would need to write for every external function call a declaration
	\item Because until now we only compiled based on function declarations there are still many possible errors ahead
	\item Header files should normally always have a matching source file and only contain function declarations
\end{itemize}
\end{frame}

\begin{frame}{C++ Compilation $\gg$ Link all objects}
\begin{itemize}
	\item The next step is to create templates, because templates are only blueprints on which basis the compiler writes functions (If you want to know more about templates which are even more powerful then for example Java Generics look here: \url{http://www.cplusplus.com/doc/oldtutorial/templates/})
	\item Then when we have object files that contain in total   all function definitions the linker tries to combine everything
	\item If in this step a function/symbol definition is not found because there was only a declaration you will experience a \textbf{Linker Error} which is easy to resolve by providing the missing definition to an already existing definition in a source file
	\item If everything \textit{links} the linker will provide you the executable (\mintinline{c}{*.exe})
\end{itemize}
\end{frame}

\begin{frame}{C++ Compilation $\gg$ A trivial C example}
\begin{columns}
	\begin{column}{0.675\textwidth}
		\begin{enumerate}
			\item We can build it with:
			\inputminted[bgcolor=lightGreyCustom,fontsize=\scriptsize]{sh}{./resources/build_trivial_c_program.sh}
		\end{enumerate}
	\end{column}
	\begin{column}{0.37\textwidth}  %%<--- here
		\begin{enumerate}\setcounter{enumi}{2}
			\item \mintinline{sh}{add.h}
			\inputminted[bgcolor=lightGreyCustom,fontsize=\scriptsize]{c}{./resources/add.h}
		\end{enumerate}
	\end{column}
\end{columns}
\vspace{-3mm}
\begin{columns}
	\begin{column}{0.4\textwidth}
		\begin{enumerate}\setcounter{enumi}{1}
			\item \mintinline{sh}{add.c}
			\inputminted[bgcolor=lightGreyCustom,fontsize=\scriptsize]{c}{./resources/add.c}
		\end{enumerate}
	\end{column}
	\begin{column}{0.6\textwidth}  %%<--- here
		\begin{enumerate}\setcounter{enumi}{3}
			\item \mintinline{sh}{simple.c} (contains main function)
			\inputminted[bgcolor=lightGreyCustom,fontsize=\scriptsize]{c}{./resources/simple.c}
		\end{enumerate}
	\end{column}
\end{columns}
\end{frame}

\begin{frame}{C++ Compilation $\gg$ \mintinline{sh}{ifndef}, \mintinline{sh}{pragma once}}
\begin{itemize}
	\item One more thing about header files and the linking process
	\item A linker error is not only occurring when to a declaration no definition exists, but also if there is more than one
	\item But thanks to the preprocessor the header files that link to each other will all be copied into one file for each compilation which means a header file could be included more than once (which could even in the compilation because of two declarations create an error)
	\item Because of this we use guards so that every header is only included once even when referenced more than once (they are the same):
\end{itemize}
\begin{columns}
	\begin{column}{0.5\textwidth}
		\begin{enumerate}
			\item The \mintinline{sh}{ifndef} guard
			\inputminted[bgcolor=lightGreyCustom,fontsize=\scriptsize]{cpp}{./resources/ifndef_example.h}
		\end{enumerate}
	\end{column}
	\begin{column}{0.5\textwidth}
		\begin{enumerate}\setcounter{enumi}{1}
			\item Or the modern variant: \mintinline{sh}{pragma once}
			\inputminted[bgcolor=lightGreyCustom,fontsize=\scriptsize]{cpp}{./resources/pragma_once_example.h}
		\end{enumerate}
	\end{column}
\end{columns}
\end{frame}

\section{CMake}

\begin{frame}{CMake $\gg$ What is CMake?}
\begin{itemize}
	\item CMake is a free open source cross platform program to help with the build process of software like C++ projects
	\item Although there are a lot more and better tools (like for example Premake) out there here is a short overview about the CMAKE pipeline
	\item The tutorial is based on a tutorial from John Lamp from 2013 \url{https://www.johnlamp.net/cmake-tutorial-1-getting-started.html}
	\item It goes without saying that you first need to install CMAKE before continuing the tutorial
\end{itemize}
\end{frame}

\begin{frame}{CMake $\gg$ ToDo List example I}
\begin{itemize}
	\item CMake has project files like Make has Makefiles
	\item \mintinline{sh}{CMakeLists.txt} describes your project to CMake and affect its output, they should be fairly simpler to read and edit compared to Makefiles
	\inputminted[bgcolor=lightGreyCustom,fontsize=\scriptsize]{cmake}{./resources/CMake_ToDo_List/CMakeLists.txt}
	\item The \mintinline{sh}{project(name)} command names your project (If the project name contains spaces it must be surrounded by quotes)
	\item The \mintinline{sh}{add_executable(target sources...)} command tells CMake you want to create an executable and adds it as target (The first arguemtn ist the name of the executable and the rest are the source files)
\end{itemize}
\end{frame}

\begin{frame}{CMake $\gg$ ToDo List example II}
\begin{columns}
	\begin{column}{0.675\textwidth}
		\begin{itemize}
			\item \mintinline{sh}{main.cpp}
			\inputminted[bgcolor=lightGreyCustom,fontsize=\scriptsize]{cpp}{./resources/CMake_ToDo_List/main.cpp}
			\item \mintinline{sh}{todo.cpp}
			\inputminted[bgcolor=lightGreyCustom,fontsize=\scriptsize]{cpp}{./resources/CMake_ToDo_List/todo.cpp}
		\end{itemize}
	\end{column}
	\begin{column}{0.325\textwidth}
		\begin{itemize}
			\item \mintinline{sh}{todo.h}
			\inputminted[bgcolor=lightGreyCustom,fontsize=\scriptsize]{cpp}{./resources/CMake_ToDo_List/todo.h}
		\end{itemize}
	\end{column}
\end{columns}
\end{frame}

\begin{frame}{CMake $\gg$ ToDo List example III}
\begin{itemize}
	\item Although CMake can do in- and out-of-source builds even  its documentation suggests to do the latter because it's much easier to have the build separated in another directory than in between the source files
	\inputminted[bgcolor=lightGreyCustom,fontsize=\scriptsize]{sh}{./resources/CMake_ToDo_List/createCmakeMakefile.sh}
\end{itemize}
\end{frame}

\section{Pointers in C++}

\begin{frame}{Pointers in C++ $\gg$ What is a pointer?}
\begin{itemize}
	\item An integer that stores a memory address
	\item You can probably better picture it with the following analogy:
	\begin{itemize}
		\item There is a straight street with a start and an end
		\item On the left side of this street there are houses with house numbers
		\item Each house has the size of a byte and it's address is called the memory address
		\item All houses together are your memory
	\end{itemize}
	\item Because everything is in memory, every variable has a memory address
	\inputminted[bgcolor=lightGreyCustom,fontsize=\scriptsize]{cpp}{./resources/memory_address.cpp}
\end{itemize}
\end{frame}

\begin{frame}{Pointers in C++ $\gg$ Purest pointer: \mintinline{cpp}{void*}}
\begin{itemize}
	\item \mintinline{cpp}{void} pointer means \textit{We do not care what the type of the pointers actual data is}
	\inputminted[bgcolor=lightGreyCustom,fontsize=\scriptsize]{cpp}{./resources/void_ptr.cpp}
\end{itemize}
\end{frame}

\begin{frame}{Pointers in C++ $\gg$ Assign memory address to pointer}
\begin{itemize}
	\item Any memory address can be assigned to a pointer, the type of the pointer changes absolutely nothing
	\inputminted[bgcolor=lightGreyCustom,fontsize=\scriptsize]{cpp}{./resources/assign_ptr.cpp}
	\item[$\Rightarrow$] This means types are only useful for manipulating data
\end{itemize}
\end{frame}

\begin{frame}{Pointers in C++ $\gg$ Change the data they point to}
\begin{itemize}
	\item Using so called dereferencing of a pointer we can change the data the pointer points to
	\item The pointer type \textbf{should} suggest to what the pointer is pointing to (If a pointer gets casted the compiler can't check it)
	\inputminted[bgcolor=lightGreyCustom,fontsize=\scriptsize]{cpp}{./resources/change_ptr_data.cpp}
\end{itemize}
\end{frame}

\begin{frame}{Pointers in C++ $\gg$ Using pointers outside the stack I}
\begin{itemize}
	\item In short: The \textbf{stack} is used for static memory allocation and the \textbf{heap} for dynamic memory allocation, both are stored in the computer's RAM
	\item In C++ the stack memory is \textit{generally} where local variables get stored/constructed, it also holds parameters passed to functions
	\item To free the stack from used memory is trivial, because the compiler already knows at compile time when a variable is never used again
	\item The stack is usually very limited compared to the heap, if you have too much memory allocated you get a stack (buffer) overflow!
	\item Heap memory is as big as the virtual memory of a computer is and a place where you can store big and more importantly dynamically sized data and then point to it
	\item On the heap is the big danger of so known memory leaks if you forget to manually free memory you allocated on it
\end{itemize}
\end{frame}

\begin{frame}{Pointers in C++ $\gg$ Using pointers outside the stack II}
\begin{itemize}
	\item Now lets for example allocate some memory on the heap and the free this memory again:
	\inputminted[bgcolor=lightGreyCustom,fontsize=\scriptsize]{cpp}{./resources/heap_ptr.cpp}
\end{itemize}
\end{frame}

\begin{frame}{Pointers in C++ $\gg$ Aren't pointer a big hazard?}
\begin{itemize}
	\item Yes they are
	\item You as a programmer need always to be fully aware to what they point to, if the data they point to was already freed (virtual memory access error), if it was changed through another pointer to the same data, if the pointer points to a different type than it says, etc.
	\item Because of this there are better pointers available that help programmers to make less errors and help with automatic memory management
	\item In the following slides we will explore smart pointers that are part of the standard library since C++11/14 and the concept of RAII
\end{itemize}
\end{frame}

\begin{frame}{Pointers in C++ $\gg$ Smart pointers: RAII}
\begin{itemize}
	\item Smart pointer were designed with compiler type safety and automatic memory management in mind
	\item The latter one is also better known as \textit{Resource Acquisition Is Initialization} programming
	\item The idea behind this idiom is to ensure that resource acquisition occurs at the same time that the object is initialized, so that all resources for the object are created and made ready in one line of code
	\item In practical terms, the main principle of RAII is to give ownership of any heap-allocated resource — for example, dynamically-allocated memory or system object handles — to a stack-allocated object whose destructor contains the code to delete or free the resource and also any associated cleanup code (A destructor of an object is like an constructor but is not called when an object is created but when it's destroyed)
\end{itemize}
\end{frame}

\begin{frame}{Pointers in C++ $\gg$ Smart pointers: Are normal pointers bad?}
\begin{itemize}
	\item Smart pointer make it possible to handle pointers without using the keyword \mintinline{cpp}{new} or \mintinline{cpp}{delete}
	\item This means smart pointers are just a wrapper around \textit{normal} pointers which is why smart pointers do not replace normal pointers - the idea is only to make the program \textit{automatically} memory safer
	\item Smart pointers are all based on pointers, but implement in an overhead special features that help in what the idea is behind them
	\item You should use smart pointers or something that works the same way, but sometimes when performance is crucial and there is no chance of confusion about ownership you are better of using old-school pointers
\end{itemize}
\end{frame}

\begin{frame}{Pointers in C++ $\gg$ Smart pointer: Class for explanation}
\begin{columns}
	\begin{column}{0.375\textwidth}
		\begin{itemize}
			\item \mintinline{sh}{ExampleClass.h}
			\inputminted[bgcolor=lightGreyCustom,fontsize=\scriptsize]{cpp}{./resources/ExampleClass.h}
		\end{itemize}
	\end{column}
	\begin{column}{0.65\textwidth}
		\begin{itemize}
			\item \mintinline{sh}{ExampleClass.cpp}
			\inputminted[bgcolor=lightGreyCustom,fontsize=\scriptsize]{cpp}{./resources/ExampleClass.cpp}
		\end{itemize}
	\end{column}
\end{columns}
\end{frame}


\begin{frame}{Pointers in C++ $\gg$ Smart pointer: 1.  \mintinline{cpp}{std::unique_ptr}}
\begin{itemize}
	\item When that pointer get out of scope it gets automatically deleted (there is no overhead using it compared to a normal pointer)
	\item You also can't copy it (\textit{only one pointer to a point in memory})
	\inputminted[bgcolor=lightGreyCustom,fontsize=\scriptsize]{cpp}{./resources/unique_ptr.cpp}
	\vspace{-9mm}
	\inputminted[bgcolor=lightGreyCustom,fontsize=\scriptsize]{sh}{./resources/build_unique_ptr.sh}
\end{itemize}
\end{frame}

\begin{frame}{Pointers in C++ $\gg$ Smart pointer: 2.  \mintinline{cpp}{std::shared_ptr} I}
\begin{itemize}
	\item Reference counting: The pointer counts in it's overhead the references to it and when the number reaches $0$ it deallocates memory (also it can be copied/shared)
	\inputminted[bgcolor=lightGreyCustom,fontsize=\scriptsize,lastline=17]{cpp}{./resources/shared_ptr.cpp}
\end{itemize}
\end{frame}

\begin{frame}{Pointers in C++ $\gg$ Smart pointer: 2.  \mintinline{cpp}{std::shared_ptr} II}
\begin{itemize}
	\inputminted[bgcolor=lightGreyCustom,fontsize=\scriptsize,firstline=18]{cpp}{./resources/shared_ptr.cpp}
	\vspace{-9mm}
	\inputminted[bgcolor=lightGreyCustom,fontsize=\scriptsize ]{sh}{./resources/build_shared_ptr.sh}
\end{itemize}
\end{frame}

\begin{frame}{Pointers in C++ $\gg$ Smart pointer: 3.  \mintinline{cpp}{std::weak_ptr} I}
\begin{itemize}
	\item Extends the \mintinline{cpp}{std::shared_ptr} but does not increase the reference count and implements a method where you can ask if its still alive (references $> 0$, Useful for something like cache interaction)
	\inputminted[bgcolor=lightGreyCustom,fontsize=\scriptsize,lastline=17]{cpp}{./resources/weak_ptr.cpp}
\end{itemize}
\end{frame}

\begin{frame}{Pointers in C++ $\gg$ Smart pointer: 3.  \mintinline{cpp}{std::weak_ptr} II}
\begin{itemize}
	\inputminted[bgcolor=lightGreyCustom,fontsize=\scriptsize,firstline=18]{cpp}{./resources/weak_ptr.cpp}
	\vspace{-9mm}
	\inputminted[bgcolor=lightGreyCustom,fontsize=\scriptsize ]{sh}{./resources/build_weak_ptr.sh}
\end{itemize}
\end{frame}

\section{Templates in C++}

\begin{frame}{Templates in C++ $\gg$ What is a template?}
\begin{itemize}
	\item Templates are somewhat comparable to Java Generics but way more powerful, because they are not limited by the type system and evaluated at compile time - they also allow custom type definitions even with the basic types
	\item Its like a macro, where you can define a blueprint of something and then the compiler will based on these rules write code for you (which is good as long as you don't need to debug it)
	\item When writing a function/class with template integration it therefore enables you to call/create the same function/object with different parameters each time you call/create the function/object
	\item The whole thing is like \textit{Meta programming}
	\item Keep in mind there are still rules which can produce many errors (even the compiler can't just find a way to print out a random custom made class object)
\end{itemize}
\end{frame}

\begin{frame}{Templates in C++ $\gg$ Example: Print function template}
\begin{itemize}
	\item In this example we want to print something but not write for every type a function that overloads itself (this would be mayor and unnecessary code duplication)
	\inputminted[bgcolor=lightGreyCustom,fontsize=\scriptsize ]{cpp}{./resources/templates/print_function_template_example.cpp}
\end{itemize}
\end{frame}

\begin{frame}{Templates in C++ $\gg$ Example: Array class template}
\begin{itemize}
	\item In this example we want to have arrays with different size and type but without evaluating it at runtime
	\inputminted[bgcolor=lightGreyCustom,fontsize=\scriptsize ]{cpp}{./resources/templates/array_class_template_example.cpp}
\end{itemize}
\end{frame}

\section{Operator overloading in C++}

\begin{frame}{Operator overloading in C++ $\gg$ Example I}
\begin{itemize}
	\item Because C++ gives you as much power as it can you can even overload operators and either making you program better or making everything worse, here is a simple example:
	\inputminted[bgcolor=lightGreyCustom,fontsize=\scriptsize,lastline=18]{cpp}{./resources/operator_overloading.cpp}
\end{itemize}
\end{frame}

\begin{frame}{Operator overloading in C++ $\gg$ Example II}
\begin{itemize}
	\inputminted[bgcolor=lightGreyCustom,fontsize=\scriptsize,firstline=19]{cpp}{./resources/operator_overloading.cpp}
\end{itemize}
\end{frame}

\begin{frame}{Issues/Ideas + Do you need to know anything else?}
	\begin{itemize}
		\item If there are any issues (things that are wrong/missing) tell me and I update the file!
	\end{itemize}
\end{frame}

\begin{frame}{Links}
	\begin{itemize}
		\item\urlCustom{github.com/green7ea/cpp-compilation}
		\item\urlCustom{www.learncpp.com}
		\item\urlCustom{en.wikipedia.org/wiki/C++}
		
	\end{itemize}
	\begin{center}\doclicenseThis\end{center}
\end{frame}

\end{document}
